\begin{abstract}
A Go engine normally gives every move the same search budget. We ask how much stronger a frozen engine becomes from two changes that leave its network and its search untouched: deciding move by move how long to search, and never searching the same position twice. Mikiri wraps KataGo with a short ladder of visit budgets, a learned rule that stops once the winrate still at stake no longer justifies the next rung, and a memory keyed on symmetry-canonical positions. On 6{,}000 positions, each searched at seven budgets and scored by independent forced searches, the stopping rule matches uniform search at 1.7 to 2.1 times the visits, ahead of ten published stopping rules re-implemented on the same engine. In 19$\times$19 play against KataGo at 200 visits per move, Mikiri at the same mean visits scores \longScore\% over \longGames{} games with a six-rung ladder ($\longElo$ Elo, 95\% interval $\longEloLo$ to $\longEloHi$) and \mainScore\% over \mainGames{} games with a four-rung one ($\mainElo$ Elo, $\mainEloLo$ to $\mainEloHi$). Doubling KataGo's own budget is worth $\unidoubleElo$. Held to \strictVisits{} visits per move, so that it requests fewer visits than its opponent even when every restarted search is counted in full, Mikiri scores \strictScore\% ($\strictElo$ Elo, $\strictEloLo$ to $\strictEloHi$). Exact positions recur often: a memory keyed on canonical positions serves \memonlyHitRate\% of moves in a 600-game series with no loss of strength, though the network evaluations it saves were already cheap, because the engine caches them. Approximate retrieval, the idea the first version of this project was built on, does not transfer: a similar stored position points to the right move 5.5\% of the time, against 43\% for the policy network's own second choice. We also show that the match records reported for that first version were produced by a protocol error, and we release the corrected testbed, dataset, models, and game records.
\end{abstract}
